% registry: SINES
\methodhead{SINES}{Sines (1959): deviatoric amplitude and mean hydrostatic stress}\label{sheet:SINES}

\block{Source} G.~Sines, in G.~Sines, J.~L.~Waisman (eds.), \emph{Metal Fatigue}, McGraw-Hill, New
York (1959) \cite{sines1959}. Implemented in the form restated and applied by Papadopoulos et al.,
\emph{Int.\ J.\ Fatigue} 19 (1997) 219--235 \cite{papadopoulos1997}, Eqs.~(62)--(64), p.~9; Table~1.

\block{Idea and scope} Stress-invariant fatigue-limit criterion, the oldest of the invariant criteria
in common use. As in Crossland (Sheet~\ref{sheet:CROSSLAND}) the amplitude of $\sqrt{J_2}$ drives
fatigue, but only the \emph{mean} hydrostatic stress shifts the limit; the hydrostatic amplitude has
no effect. Proposed for the high-cycle fatigue limit of metals with mean stresses.

\block{Inputs}
\begin{center}\small
\begin{tabular}{@{}llp{100mm}@{}}
\toprule
Symbol & Unit & Meaning\\\midrule
$\fm$, $\tm$ & MPa & fully reversed tension (bending) and torsion fatigue limits\\
$f_0$ & MPa & repeated ($R=0$) bending limit in Eq.~\eqref{eq:sines-63}; with the Goodman estimate below it is the amplitude (equal to the mean) of the $R=0$ cycle\\
$\sigma_u=S_u$ & MPa & ultimate tensile strength (written $\sigma_f$ in the paper's Eq.~64)\\
$\sqrt{J_{2,a}}$ & MPa & amplitude of $\sqrt{J_2}$, the enclosing radius defined in Sheet~\ref{sheet:CROSSLAND}\\
$\sigma_{H,m}$ & MPa & mean hydrostatic stress, $\tr\Sm/3$\\
$\kappa$, $\lambda$ & --, MPa & constants of the criterion (this $\kappa$ is not the ratio of Section~\ref{sec:notation})\\
$\mathrm{DP}$ & -- & fatigue index, Eq.~\eqref{eq:sines-index}\\
\bottomrule
\end{tabular}
\end{center}

\block{Formulation} Papadopoulos et al.\ write Sines' criterion as (Eq.~62)
\begin{equation}
\sqrt{J_{2,a}} + \kappa\,\sigma_{H,m} \le \lambda .
\label{eq:sines-main}
\end{equation}
Sines' own calibration (Eq.~63) uses the fully reversed and repeated bending limits,
\begin{equation}
\kappa = \frac{3\,\tm}{f_0} - \sqrt3 , \qquad \lambda = \tm .
\label{eq:sines-63}
\end{equation}
With the Goodman estimate $f_0=\fm/(1+\fm/\sigma_u)$ and the torsion-to-bending ratio implied by the
criterion, $\tm=\fm/\sqrt3$, this becomes the form evaluated in the paper's Table~1 (Eq.~64):
\begin{equation}
\sqrt{J_{2,a}} + \kappa\,\sigma_{H,m} \le \tm , \qquad \kappa = \frac{\sqrt3\,\fm}{\sigma_u} .
\label{eq:sines-64}
\end{equation}
The repository returns the index
\begin{equation}
\mathrm{DP} = \frac{\sqrt{J_{2,a}} + \kappa\,\sigma_{H,m}}{\tm} , \qquad \kappa = \frac{\sqrt3\,\fm}{S_u} .
\label{eq:sines-index}
\end{equation}

\block{Constants and calibration} Paper: $\kappa$ from the reversed and repeated bending limits
(Eq.~63). Repository: the closed form of Eq.~\eqref{eq:sines-64} from the workbook $\fm$, $\tm$ and
$S_u$; a measured $f_0$ (\texttt{Sig0}) is not used. Only the master-curve constants are fitted, on
the uniaxial $R=-1$ design tests.

\block{Life} No native life law. The index is the damage parameter of the shared master curve
(Section~\ref{sec:mastercurve}), a repository extension of a fatigue-limit criterion.

\block{Repository choices}
\begin{itemize}
\item The printed Eq.~(64) coefficient $\kappa=\sqrt3\,\fm/S_u$, independent of $\tm$, as in the
  paper's Table~1. Fully reversed tension at $\fm$ therefore gives the index $\fm/(\sqrt3\,\tm)$,
  which is 1 only when $\tm=\fm/\sqrt3$. Eq.~(63) with the Goodman $f_0$ would miss the paper's
  mean-stress rows by up to 11 points.
\item $\sqrt{J_{2,a}}$ is the exact enclosing radius of the single-harmonic deviatoric ellipse, as in
  Crossland.
\item The index is the damage parameter of the master curve, for finite-life comparisons.
\item Case-level model: no plane search and no \texttt{\_ext} variant.
\end{itemize}

\block{Worked example} Synthetic material $\fm=300$, $\tm=200$, $S_u=800\MPa$: $\kappa=0.6495$. Torsion
at $\tm$ gives index~1; tension at $\fm$ gives $0.8660$ ($=\fm/(\sqrt3\,\tm)$). Tension
$\sigma_a=200\MPa$ with mean $\sigma_m=150\MPa$: index $0.7397$, equal to
$(\sigma_a/\sqrt3+\kappa\,\sigma_m/3)/\tm$.

\block{Verification} Papadopoulos et al.\ (1997) Table~1 is reproduced for $\delta=0^\circ$,
$90^\circ$ and $180^\circ$ with the printed Eq.~(64). The rows at intermediate phase shifts are
strict xfails: the table used $\sin\delta$ for the $\sin^2\delta$ of Eq.~(44). Rows~14 and~15
(42CrMo4) are 0.3 points off and kept as strict xfails (unresolved).
